% --- CORRECCION: Forzar interpretacion raw de la entrada ---
\UseRawInputEncoding
\documentclass[aspectratio=169, 10pt, xcolor=table]{beamer}

% --- Configuracion del Tema ---
\usetheme{Madrid}
\usecolortheme{dolphin}

% --- Paquetes necesarios ---
\usepackage[utf8]{inputenc}
\usepackage[T1]{fontenc}
\usepackage[spanish,es-noquoting]{babel}
\usepackage{amsmath, amssymb}
\usepackage{mathtools}
\usepackage{booktabs}
\usepackage{graphicx}
\usepackage{listings}
\usepackage{tikz}
\usetikzlibrary{arrows.meta, positioning, shapes.geometric, babel}
\usepackage{etoolbox}
\usepackage{upquote}

% --- Estilo para codigo SQL ---
\definecolor{codegreen}{rgb}{0,0.6,0}
\definecolor{codegray}{rgb}{0.5,0.5,0.5}
\definecolor{codepurple}{rgb}{0.58,0,0.82}
\definecolor{backcolour}{rgb}{0.95,0.95,0.92}
\lstset{
    language=SQL,
    backgroundcolor=\color{backcolour},
    commentstyle=\color{codegreen},
    keywordstyle=\color{blue},
    numberstyle=\tiny\color{codegray},
    stringstyle=\color{codepurple},
    basicstyle=\ttfamily\scriptsize,
    breaklines=true,
    frame=single,
    showstringspaces=false,
    literate={á}{{\'a}}1 {é}{{\'e}}1 {í}{{\'i}}1 {ó}{{\'o}}1 {ú}{{\'u}}1 {ñ}{{\~n}}1 {ü}{{\"u}}1
}

% --- Pie de pagina personalizado ---
\makeatletter
\setbeamertemplate{footline}{
  \leavevmode%
  \hbox{%
  \begin{beamercolorbox}[wd=.333333\paperwidth,ht=2.25ex,dp=1ex,center]{author in head/foot}%
    \usebeamerfont{author in head/foot}\insertshortauthor
  \end{beamercolorbox}%
  \begin{beamercolorbox}[wd=.333333\paperwidth,ht=2.25ex,dp=1ex,center]{title in head/foot}%
    \usebeamerfont{title in head/foot}Sesión 5
  \end{beamercolorbox}%
  \begin{beamercolorbox}[wd=.333333\paperwidth,ht=2.25ex,dp=1ex,right]{date in head/foot}%
    \usebeamerfont{date in head/foot}BBDD \hfill \insertframenumber/\inserttotalframenumber \hspace*{0.3cm}
  \end{beamercolorbox}}%
  \vskip0pt%
}
\makeatother

% --- Datos de la portada ---
\title[SESIÓN 5 BBDD]{\textbf{Sesión 5}}
\subtitle{Manipulación de Datos y Lógica de Conjuntos}
\author[Prof. C. Sánchez]{Carlos César Sánchez Coronel}
\institute{\textbf{BASES DE DATOS}}
\date{}

\setbeamertemplate{navigation symbols}{}

\AtBeginSection[]{
  \begin{frame}
    \frametitle{Contenido de la sección}
    \tableofcontents[currentsection]
  \end{frame}
}

\begin{document}

\begin{frame}
    \titlepage
\end{frame}

\begin{frame}{Contenido general}
    \tableofcontents
\end{frame}

% ============================================================================
% SECCIÓN 1: INTRODUCCIÓN A LA MANIPULACIÓN DE DATOS Y SUBLENGUAJES SQL
% ============================================================================
\section{Introducción y Sublenguajes SQL}

\begin{frame}{Sublenguajes SQL}
    \begin{block}{Clasificación}
        \begin{itemize}
            \item \textbf{DDL} (Data Definition Language): CREATE, ALTER, DROP, TRUNCATE.
            \item \textbf{DML} (Data Manipulation Language): SELECT, INSERT, UPDATE, DELETE, MERGE.
            \item \textbf{DCL} (Data Control Language): GRANT, REVOKE.
            \item \textbf{TCL} (Transaction Control Language): BEGIN, COMMIT, ROLLBACK, SAVEPOINT.
        \end{itemize}
    \end{block}
\end{frame}

\begin{frame}[fragile]{DDL: Creación y modificación de estructuras}
    \begin{exampleblock}{CREATE TABLE}
        \begin{lstlisting}
CREATE TABLE clientes (
    id SERIAL PRIMARY KEY,
    nombre VARCHAR(100) NOT NULL,
    email VARCHAR(100) UNIQUE,
    ciudad VARCHAR(50),
    fecha_registro DATE DEFAULT CURRENT_DATE
);
        \end{lstlisting}
    \end{exampleblock}
    \begin{exampleblock}{ALTER y DROP}
        \begin{lstlisting}
ALTER TABLE clientes ADD COLUMN telefono VARCHAR(20);
DROP TABLE clientes; -- Cuidado!
        \end{lstlisting}
    \end{exampleblock}
\end{frame}

\begin{frame}[fragile]{DML: Manipulación de datos}
    \begin{columns}[t]
        \column{0.5\textwidth}
        \textbf{INSERT}
        \begin{lstlisting}
INSERT INTO clientes (nombre, email, ciudad)
VALUES ('Carlos Ruiz', 'carlos@mail.com', 'Valencia');
        \end{lstlisting}
        \textbf{UPDATE}
        \begin{lstlisting}
UPDATE clientes SET ciudad = 'Sevilla'
WHERE id = 5;
        \end{lstlisting}
        \column{0.5\textwidth}
        \textbf{DELETE}
        \begin{lstlisting}
DELETE FROM clientes WHERE email IS NULL;
        \end{lstlisting}
        \textbf{MERGE (UPSERT en PostgreSQL)}
        \begin{lstlisting}
INSERT INTO productos (id, nombre, precio)
VALUES (1, 'Laptop', 1200)
ON CONFLICT (id) DO UPDATE SET
    nombre = EXCLUDED.nombre,
    precio = EXCLUDED.precio;
        \end{lstlisting}
    \end{columns}
\end{frame}

\begin{frame}[fragile]{DCL y TCL}
    \begin{block}{DCL: Control de permisos}
        \begin{lstlisting}
GRANT SELECT ON clientes TO analista;
REVOKE DELETE ON clientes FROM analista;
        \end{lstlisting}
    \end{block}
    \begin{block}{TCL: Transacciones}
        \begin{lstlisting}
BEGIN;
UPDATE cuentas SET saldo = saldo - 100 WHERE id = 1;
UPDATE cuentas SET saldo = saldo + 100 WHERE id = 2;
COMMIT; -- o ROLLBACK si algo falla
        \end{lstlisting}
    \end{block}
\end{frame}

\begin{frame}{Propiedades ACID}
    \begin{description}
        \item[Atomicidad] Todo o nada.
        \item[Consistencia] Integridad de datos.
        \item[Aislamiento] Transacciones concurrentes no interfieren.
        \item[Durabilidad] Los cambios persisten tras fallos.
    \end{description}
\end{frame}

% ============================================================================
% SECCIÓN 2: PROCESO INTERNO DE UNA CONSULTA
% ============================================================================
\section{Proceso Interno de una Consulta}

\begin{frame}{Orden Lógico de Ejecución}
    \begin{enumerate}
        \item FROM y JOINs
        \item WHERE
        \item GROUP BY
        \item HAVING
        \item SELECT (y funciones de ventana)
        \item ORDER BY
        \item LIMIT / OFFSET
    \end{enumerate}
    \begin{alertblock}{Importante}
        Los alias de columna definidos en SELECT no pueden usarse en WHERE (porque WHERE se ejecuta antes).
    \end{alertblock}
\end{frame}

\begin{frame}[fragile]{El Camino Físico}
    \begin{enumerate}
        \item \textbf{Parsing}: Verificación sintáctica y semántica.
        \item \textbf{Optimización}: Generación de planes de ejecución (basado en costos).
        \item \textbf{Ejecución}: Interacción con buffer pool y disco.
        \item \textbf{Retorno}: Envío de resultados al cliente.
    \end{enumerate}
    \begin{block}{EXPLAIN (ANALYZE)}
        \begin{verbatim}
EXPLAIN (ANALYZE, BUFFERS) 
SELECT * FROM clientes WHERE ciudad = 'Madrid';
        \end{verbatim}
    \end{block}
\end{frame}

% ============================================================================
% SECCIÓN 3: JOINS
% ============================================================================
\section{Joins: Combinando Tablas}

\begin{frame}{Tipos de JOIN}
    \begin{columns}[t]
        \column{0.5\textwidth}
        \textbf{INNER JOIN}
        \begin{itemize}
            \item Solo filas con coincidencia en ambas.
        \end{itemize}
        \textbf{LEFT JOIN}
        \begin{itemize}
            \item Todas las filas de la izquierda + coincidencias.
        \end{itemize}
        \column{0.5\textwidth}
        \textbf{RIGHT JOIN}
        \begin{itemize}
            \item Todas las filas de la derecha + coincidencias.
        \end{itemize}
        \textbf{FULL OUTER JOIN}
        \begin{itemize}
            \item Todas las filas de ambas, con NULL donde no hay.
        \end{itemize}
    \end{columns}
\end{frame}

\begin{frame}[fragile]{Ejemplos de JOIN}
    \begin{block}{INNER JOIN}
        \begin{lstlisting}
SELECT p.id, c.nombre, p.fecha
FROM pedidos p
INNER JOIN clientes c ON p.cliente_id = c.id;
        \end{lstlisting}
    \end{block}
    \begin{block}{LEFT JOIN (clientes con o sin pedidos)}
        \begin{lstlisting}
SELECT c.nombre, COUNT(p.id) AS num_pedidos
FROM clientes c
LEFT JOIN pedidos p ON c.id = p.cliente_id
GROUP BY c.id, c.nombre;
        \end{lstlisting}
    \end{block}
\end{frame}

\begin{frame}[fragile]{Self-Join y CROSS JOIN}
    \begin{block}{Self-Join: Empleados con su jefe}
        \begin{lstlisting}
SELECT e1.nombre AS empleado, e2.nombre AS jefe
FROM empleados e1
LEFT JOIN empleados e2 ON e1.jefe_id = e2.id;
        \end{lstlisting}
    \end{block}
    \begin{block}{CROSS JOIN: Producto cartesiano}
        \begin{lstlisting}
SELECT * FROM colores CROSS JOIN tallas;
        \end{lstlisting}
    \end{block}
\end{frame}

\begin{frame}{Algoritmos Internos de JOIN}
    \begin{description}
        \item[Nested Loop:] Una tabla pequeña o con índice; va fila por fila.
        \item[Hash Join:] Tablas grandes sin índice; construye una tabla hash.
        \item[Merge Join:] Tablas ordenadas; fusiona ambas.
    \end{description}
\end{frame}

% ============================================================================
% SECCIÓN 4: SUBCONSULTAS
% ============================================================================
\section{Subconsultas}

\begin{frame}[fragile]{Tipos de Subconsultas}
    \begin{block}{Escalar (un solo valor)}
        \begin{lstlisting}
SELECT nombre, (SELECT AVG(precio) FROM productos) AS promedio
FROM productos;
        \end{lstlisting}
    \end{block}
    \begin{block}{Múltiples filas con IN}
        \begin{lstlisting}
SELECT nombre FROM clientes
WHERE id IN (SELECT cliente_id FROM pedidos WHERE monto > 1000);
        \end{lstlisting}
    \end{block}
\end{frame}

\begin{frame}[fragile]{Subconsultas Correlacionadas}
    \begin{block}{Empleados con salario mayor al promedio de su departamento}
        \begin{lstlisting}
SELECT e.nombre, e.salario
FROM empleados e
WHERE e.salario > (SELECT AVG(salario)
                   FROM empleados
                   WHERE departamento_id = e.departamento_id);
        \end{lstlisting}
    \end{block}
    \begin{alertblock}{Rendimiento}
        Se evalúan para cada fila exterior; pueden ser lentas en tablas grandes.
    \end{alertblock}
\end{frame}

\begin{frame}[fragile]{Subconsultas en FROM}
    \begin{block}{Tabla derivada}
        \begin{lstlisting}
SELECT categoria, AVG(precio) AS precio_medio
FROM (SELECT * FROM productos WHERE activo = true) AS productos_activos
GROUP BY categoria;
        \end{lstlisting}
    \end{block}
\end{frame}

% ============================================================================
% SECCIÓN 5: COMMON TABLE EXPRESSIONS (CTEs)
% ============================================================================
\section{CTEs}

\begin{frame}[fragile]{CTEs Básicas}
    \begin{block}{WITH ...}
        \begin{lstlisting}
WITH ventas_cliente AS (
    SELECT c.id, c.nombre, SUM(dp.cantidad * pr.precio) AS total
    FROM clientes c
    LEFT JOIN pedidos p ON c.id = p.cliente_id
    LEFT JOIN detalle_pedido dp ON p.id = dp.pedido_id
    LEFT JOIN productos pr ON dp.producto_id = pr.id
    GROUP BY c.id, c.nombre
)
SELECT * FROM ventas_cliente WHERE total > 1000;
        \end{lstlisting}
    \end{block}
\end{frame}

\begin{frame}[fragile]{CTEs Múltiples y Recursivas}
    \begin{block}{Múltiples CTEs}
        \begin{lstlisting}
WITH
    ventas AS (SELECT ...),
    clientes_vip AS (SELECT ...)
SELECT ...
        \end{lstlisting}
    \end{block}
    \begin{block}{Recursiva (organigrama)}
        \begin{lstlisting}
WITH RECURSIVE organigrama AS (
    SELECT id, nombre, jefe_id, 1 AS nivel
    FROM empleados WHERE jefe_id IS NULL
    UNION ALL
    SELECT e.id, e.nombre, e.jefe_id, o.nivel + 1
    FROM empleados e JOIN organigrama o ON e.jefe_id = o.id
)
SELECT * FROM organigrama ORDER BY nivel;
        \end{lstlisting}
    \end{block}
\end{frame}

% ============================================================================
% SECCIÓN 6: BUENAS PRÁCTICAS Y OPTIMIZACIÓN
% ============================================================================
\section{Buenas Prácticas y Optimización}

\begin{frame}{Consejos de Optimización}
    \begin{itemize}
        \item Preferir \texttt{EXISTS} a \texttt{IN} para subconsultas grandes.
        \item Evitar funciones en columnas indexadas en \texttt{WHERE}.
        \item Usar \texttt{UNION ALL} en lugar de \texttt{UNION} si no se necesitan duplicados.
        \item Evitar \texttt{SELECT *} (especificar columnas).
        \item Cuidado con conversiones implícitas.
        \item Indexar columnas usadas en \texttt{JOIN}, \texttt{WHERE}, \texttt{ORDER BY}.
        \item Analizar siempre con \texttt{EXPLAIN}.
    \end{itemize}
\end{frame}

\begin{frame}[fragile]{Ejemplo de Optimización}
    \begin{block}{Consulta ineficiente}
        \begin{verbatim}
SELECT * FROM pedidos WHERE EXTRACT(YEAR FROM fecha) = 2025;
        \end{verbatim}
    \end{block}
    \begin{block}{Versión optimizada (aprovecha índice en fecha)}
        \begin{verbatim}
SELECT * FROM pedidos 
WHERE fecha >= '2025-01-01' AND fecha < '2026-01-01';
        \end{verbatim}
    \end{block}
\end{frame}

% ============================================================================
% SECCIÓN 7: HARDWARE Y NOSQL
% ============================================================================
\section{Hardware y Breve Mirada a NoSQL}

\begin{frame}{Consideraciones de Hardware}
    \begin{itemize}
        \item \textbf{Buffer Pool}: Memoria caché de datos e índices (ej. \texttt{shared\_buffers} en PostgreSQL).
        \item \textbf{SSD vs HDD}: SSD mejora lecturas aleatorias, crítico para índices.
        \item \textbf{Parallel Query}: Uso de múltiples CPUs para consultas pesadas.
    \end{itemize}
\end{frame}

\begin{frame}{Joins en NoSQL}
    \begin{itemize}
        \item \textbf{MongoDB}: \texttt{\$lookup} en pipeline de agregación (left outer join).
        \item \textbf{Cassandra}: No soporta joins; modelado desnormalizado.
        \item \textbf{Redis}: No tiene joins; se maneja en aplicación.
        \item \textbf{Neo4j}: Relaciones nativas (navegación en grafos).
    \end{itemize}
\end{frame}

% ============================================================================
% SECCIÓN 8: EJERCICIOS RESUELTOS (DEL SOLUCIONARIO)
% ============================================================================
\section{Ejercicios Resueltos}

\begin{frame}[fragile]{Ejercicio 1: Consultas Básicas}
    \begin{block}{Enunciado}
        Tabla \texttt{empleados} (id, nombre, salario, departamento\_id). Insertar tres empleados, actualizar salario de uno, eliminar empleados de un departamento.
    \end{block}
    \begin{block}{Solución}
        \begin{lstlisting}
INSERT INTO empleados (nombre, salario, departamento_id) VALUES
('Juan Perez', 3000, 1),
('Ana Garcia', 3500, 2),
('Luis Fernandez', 3200, 1);

UPDATE empleados SET salario = 3300 WHERE nombre = 'Juan Perez';

DELETE FROM empleados WHERE departamento_id = 2;
        \end{lstlisting}
    \end{block}
\end{frame}

\begin{frame}[fragile]{Ejercicio 2: INNER JOIN y LEFT JOIN}
    \begin{block}{Enunciado}
        Con las tablas \texttt{clientes}, \texttt{pedidos}, \texttt{detalle\_pedido}, \texttt{productos}, obtener:
        \begin{enumerate}
            \item Pedidos con nombre de cliente y fecha.
            \item Clientes con número de pedidos (0 si no tienen).
            \item Total gastado por cliente.
        \end{enumerate}
    \end{block}
\end{frame}

\begin{frame}[fragile]{Solución Ejercicio 2}
    \begin{lstlisting}
-- 1. Pedidos con cliente
SELECT p.id, c.nombre, p.fecha
FROM pedidos p INNER JOIN clientes c ON p.cliente_id = c.id;

-- 2. Clientes con numero de pedidos
SELECT c.id, c.nombre, COUNT(p.id) AS num_pedidos
FROM clientes c LEFT JOIN pedidos p ON c.id = p.cliente_id
GROUP BY c.id, c.nombre;

-- 3. Total gastado por cliente
SELECT c.id, c.nombre, COALESCE(SUM(dp.cantidad * pr.precio), 0) AS total
FROM clientes c
LEFT JOIN pedidos p ON c.id = p.cliente_id
LEFT JOIN detalle_pedido dp ON p.id = dp.pedido_id
LEFT JOIN productos pr ON dp.producto_id = pr.id
GROUP BY c.id, c.nombre;
    \end{lstlisting}
\end{frame}

\begin{frame}[fragile]{Ejercicio 3: Subconsultas}
    \begin{block}{Enunciado}
        Productos cuyo precio es superior al precio medio de su categoría (subconsulta correlacionada).
    \end{block}
    \begin{block}{Solución}
        \begin{lstlisting}
SELECT p.nombre, p.precio
FROM productos p
WHERE p.precio > (SELECT AVG(precio)
                  FROM productos
                  WHERE categoria = p.categoria);
        \end{lstlisting}
    \end{block}
\end{frame}

\begin{frame}[fragile]{Ejercicio 4: CTE}
    \begin{block}{Enunciado}
        Obtener empleados que ganan más que el promedio de su departamento, mostrando nombre, salario y promedio del departamento.
    \end{block}
    \begin{block}{Solución}
        \begin{lstlisting}
WITH promedio_depto AS (
    SELECT departamento_id, AVG(salario) AS salario_prom
    FROM empleados GROUP BY departamento_id
)
SELECT e.nombre, e.salario, pd.salario_prom
FROM empleados e
JOIN promedio_depto pd ON e.departamento_id = pd.departamento_id
WHERE e.salario > pd.salario_prom;
        \end{lstlisting}
    \end{block}
\end{frame}

\begin{frame}[fragile]{Ejercicio 5: Optimización}
    \begin{block}{Consulta lenta}
        \begin{verbatim}
SELECT * FROM pedidos WHERE EXTRACT(YEAR FROM fecha) = 2025;
        \end{verbatim}
    \end{block}
    \begin{block}{Optimización}
        \begin{itemize}
            \item Reescribir: \texttt{WHERE fecha >= '2025-01-01' AND fecha < '2026-01-01'}
            \item Crear índice en \texttt{fecha}.
            \item Verificar con \texttt{EXPLAIN ANALYZE}.
        \end{itemize}
    \end{block}
\end{frame}

\begin{frame}[fragile]{Ejercicio 6: CTE Recursiva}
    \begin{block}{Enunciado}
        Obtener organigrama completo con niveles desde \texttt{empleados(id, nombre, jefe\_id)}.
    \end{block}
    \begin{block}{Solución}
        \begin{lstlisting}
WITH RECURSIVE organigrama AS (
    SELECT id, nombre, jefe_id, 1 AS nivel
    FROM empleados WHERE jefe_id IS NULL
    UNION ALL
    SELECT e.id, e.nombre, e.jefe_id, o.nivel + 1
    FROM empleados e JOIN organigrama o ON e.jefe_id = o.id
)
SELECT * FROM organigrama ORDER BY nivel, nombre;
        \end{lstlisting}
    \end{block}
\end{frame}

% ============================================================================
% SECCIÓN 9: RESUMEN
% ============================================================================
\section{Resumen}

\begin{frame}{Conceptos Clave de la Sesión 5}
    \begin{itemize}
        \item \textbf{Sublenguajes SQL}: DDL, DML, DCL, TCL.
        \item \textbf{Transacciones}: ACID, COMMIT, ROLLBACK, SAVEPOINT.
        \item \textbf{Orden lógico de ejecución} y su impacto.
        \item \textbf{Joins}: INNER, LEFT, RIGHT, FULL, CROSS, self-join.
        \item \textbf{Subconsultas}: escalares, correlacionadas, en FROM.
        \item \textbf{CTEs}: simples, múltiples, recursivas.
        \item \textbf{Optimización}: índices, evitar funciones en WHERE, usar EXISTS, analizar con EXPLAIN.
        \item \textbf{Hardware}: buffer pool, SSD, parallel query.
        \item \textbf{NoSQL}: aproximación a joins (\$lookup, desnormalización).
    \end{itemize}
\end{frame}

\begin{frame}
    \centering

    
    \vspace{1cm}
    \large ¡Gracias por su atención!
\end{frame}

\end{document}